\subsection{p-根式元素}

定义1. --- 设K是一个域，E是K的一个扩张。E的一个元素x称为在K上是p-根式的，如果存在一个整数 \(m \geqslant 0\) 使得 \(x^{p^{m}} \in K\) ；这些整数中最小的一个称为x（在K上）的高度。

命题1. --- 设E是域K的一个扩张，x是K上高度为e的p-根式元素；令 \(a = x^{p^{e}}\) 。则 \(a \in K\) 且x在K上的极小多项式为 \(X^{p^{e}} - a\) 。进一步我们有 \([\mathbf{K}(x): \mathbf{K}] = p^{e}\) 。

只需证明多项式 \(X^{p^{e}} - a\) 在 K{[}X{]} 中不可约。由x的高度的定义，我们有 \(a \notin K^{p}\) ，因此该命题由以下引理得出：

引理1. --- 设K是一个域，a是K的一个元素，使得 \(a \notin \mathbf{K}^p\) 。那么对于每个整数 \(e \geqslant 0\) ，多项式 \(f(\mathbf{X}) = \mathbf{X}^{p^e} - a\) 在 \(\mathbf{K}[\mathbf{X}]\) 中不可约。

设 \(\Omega\) 为 \(\mathbf{K}\) 的一个代数闭包，并且令 \(b\) 为元素 \(a^{p^{-e}}\) （在 \(\Omega\) 中）；我们用 \(g\) 表示 \(b\) 在 \(\mathbf{K}\) 上的极小多项式。我们有 \(f(\mathbf{X}) = (\mathbf{X} - b)^{p^e}\) ，因此 \(\mathbf{K}[\mathbf{X}]\) 中每个整除 \(f\) 的不可约多项式都以 \(b\) 为根，从而等于 \(g\) 。于是存在（IV，第～13页，命题13）一个整数 \(q \geqslant 1\) 使得 \(f = g^q\) ；由于 \(q\) 整除 \(p^e\) （\(f\) 的次数），存在一个整数 \(e'\) 使得 \(0 \leqslant e' \leqslant e\) 且 \(q = p^{e'}\) 。若 \(c\) 是 \(g\) 的常数项，我们有 \(-a = c^q\) ；由于我们假设 \(a\) 不属于 \(\mathbf{K}^p\) ，我们必须有 \(q = 1\) ，即 \(f = g\) ，于是引理得证。

